% !TEX root = ../main.tex
\chapter{\ifproject%
\ifenglish Project Structure and Methodology\else โครงสร้างและขั้นตอนการทำงาน\fi
\else%
\ifenglish Project Structure\else โครงสร้างของโครงงาน\fi
\fi
}

ในบทนี้จะนำเสนอสถาปัตยกรรมและการออกแบบเชิงวิศวกรรมของระบบจำลองการลงคะแนนเสียงอิเล็กทรอนิกส์ความปลอดภัยสูง โดยประยุกต์ใช้การกระจายกุญแจลับเชิงควอนตัม (BB84) และการเข้ารหัสบัตรลงคะแนนเสียงด้วย One-Time Pad (OTP) เนื้อหาถูกจัดเรียงตามหลักการต่อจิ๊กซอว์ทางวิศวกรรม เริ่มต้นจากการทำความเข้าใจบทบาทของแต่ละฝ่ายในสถาปัตยกรรมสองช่องทาง นำไปสู่ขั้นตอนการสร้างกุญแจควอนตัม การลงคะแนนเสียง และรายละเอียดการพัฒนาซอฟต์แวร์ทั้งในส่วนของฝั่งผู้ใช้งาน (Frontend) เซิร์ฟเวอร์ประมวลผล (Backend) และระบบฐานข้อมูล (Database)

% ============================================================================
% 3.1 Overall System Architecture
% ============================================================================
\section{Overall System Architecture}
การออกแบบระบบจำลองการลงคะแนนเสียงนี้มีเป้าหมายเพื่อสร้างความปลอดภัยแบบ Information-Theoretic Security โดยแก้ปัญหาเรื่องการแจกจ่ายกุญแจสุ่มสำหรับ One-Time Pad โดยใช้ Quantum Key Distribution ผ่านการจำลองโดยใช้ควอนตัมคอมพิวเตอร์สำหรับการจำลองการสื่อสารผ่าน Quantum Channel

\subsection{Dual-Channel Communication Model and Roles}
ตามหลักการของโพรโทคอล BB84 ระบบได้รับการออกแบบให้ทำงานโดยใช้ การสื่อสารสองช่องทาง โดยมีการกำหนดบทบาทของระบบการเลือกตั้งและผู้ลงคะแนนเสียงดังนี้:
\begin{itemize}
    \item \textbf{Election Server รับบทบาทเป็นผู้ส่ง:} เซิร์ฟเวอร์การเลือกตั้งทำหน้าที่เป็นฝ่ายเตรียมสถานะควอนตัม โดยทำการสุ่มบิต สุ่มเลือกแกน (Basis) ในการส่ง
    \item \textbf{Voter Client รับบทบาทเป็นผู้รับ:} อุปกรณ์หรือเว็บแอปพลิเคชันของผู้มีสิทธิ์ลงคะแนนเสียง ทำหน้าที่เป็นฝ่ายวัดสถานะควอนตัม (Quantum Measurement) โดยสุ่มเลือกแกนวัดของตนเองเพื่ออ่านค่าคิวบิตที่ได้รับ และนำกุญแจลับสุดท้ายที่ได้ไปใช้ในการเข้ารหัสบัตรลงคะแนนเสียง
\end{itemize}

การสื่อสารระหว่าง Server และ Voter ดำเนินการผ่าน 2 ช่องทาง:
\begin{enumerate}
    \item \textbf{ช่องสัญญาณควอนตัม (Quantum Channel):} เป็นช่องทางการส่งอนุภาคโฟตอนเดี่ยวที่ใช้เป็นข้อมูลคิวบิต จาก Server ไปยัง Voter ช่องทางนี้ไม่มีการเข้ารหัสด้วยคณิตศาสตร์ แต่ได้รับการคุ้มครองด้วยกฎฟิสิกส์ควอนตัม (No-Cloning Theorem) ทำให้หากมีผู้ดักฟังแอบวัดค่าโฟตอน สถานะจะเกิดการรบกวนและทิ้งร่องรอยความผิดพลาดไว้ในระบบทันที
    \item \textbf{ช่องสัญญาณคลาสสิก (Classical Channel):} เป็นการเชื่อมต่อเครือข่ายอินเทอร์เน็ตปกติผ่านโพรโทคอล HTTP/REST API ใช้สำหรับส่งข้อมูลควบคุมที่ไม่เป็นความลับ ได้แก่ การเทียบแกนที่ใช้ในการวัด (Basis Sifting) การแลกเปลี่ยนตำแหน่งบิตตัวอย่างเพื่อคำนวณค่า QBER การส่งค่า Seed สำหรับ Toeplitz Matrix ตลอดจนการส่งบัตรลงคะแนนที่ถูกเข้ารหัสแล้ว (Ciphertext)
\end{enumerate}

% ตารางที่~\ref{tab:dual_channel_comparison} สรุปการเปรียบเทียบคุณลักษณะและหน้าที่ของช่องสัญญาณทั้งสองประเภทในระบบ

\begin{table}[htbp]
\centering
\renewcommand{\arraystretch}{1.35}
\setlength{\extrarowheight}{2.5pt}
\begin{tabularx}{\textwidth}{>{\raggedright\arraybackslash}p{3.5cm} >{\raggedright\arraybackslash}X >{\raggedright\arraybackslash}X}
\hline
\textbf{คุณลักษณะ} & \textbf{ช่องสัญญาณควอนตัม} \newline \textbf{(Quantum Channel)} & \textbf{ช่องสัญญาณคลาสสิก} \newline \textbf{(Classical Channel)} \\ \hline
\rowcolor{white}
\textbf{ข้อมูลที่ส่งผ่าน} & สถานะโพลาไรเซชันของโฟตอน (Qubits) & ข้อมูลดิจิทัลทั่วไป (ฐานวัด, บิตตัวอย่าง, Ciphertext) \\
\rowcolor[HTML]{E4E7EB}
\textbf{การคุ้มครองความปลอดภัย} & กฎฟิสิกส์ควอนตัม (No-Cloning, Collapse) & สมมติฐานว่าเปิดเผยต่อสาธารณะ ดักฟังได้อิสระ \\
\rowcolor{white}
\textbf{ผลกระทบเมื่อถูกดักฟัง} & เกิดข้อผิดพลาดในคิวบิตอย่างน้อย $25\%$ & ผู้ดักฟังทราบฐานวัด แต่ไม่ได้กุญแจลับ \\
\rowcolor[HTML]{E4E7EB}
\textbf{การนำไปใช้ในระบบ} & จำลองผ่าน Qiskit Aer / IBM Quantum & พัฒนาผ่าน RESTful API (FastAPI) บนเครือข่าย \\ \hline
\end{tabularx}
\caption{การเปรียบเทียบคุณลักษณะระหว่างช่องสัญญาณควอนตัมและช่องสัญญาณสาธารณะคลาสสิก}
\label{tab:dual_channel_comparison}
\end{table}

\subsection{Three-Tier Software Architecture}
เพื่อให้ระบบมีความยืดหยุ่น ดูแลรักษาง่าย และรองรับการขยายตัว โครงสร้างซอฟต์แวร์จึงถูกออกแบบตามสถาปัตยกรรมแบบ 3 ชั้น (Three-Tier Architecture) ดังนี้:
\begin{enumerate}
    \item \textbf{Presentation Tier (Frontend Web Client):} พัฒนาด้วย React, Vite และ TypeScript ทำหน้าที่ติดต่อกับผู้ใช้งาน โดยออกแบบส่วนติดต่อผู้ใช้งานให้สามารถแสดงผลได้ 3 มุมมอง (Multi-Perspective View) เพื่อประโยชน์ทั้งในการลงคะแนนเสียงจริงและการศึกษาเรียนรู้ ได้แก่ Voter View, Practical View และ System View
    \item \textbf{Logic Tier (Backend Application Server):} พัฒนาด้วยภาษา Python บนเฟรมเวิร์ค FastAPI ทำหน้าที่ประมวลผล ควบคุมเซสชันการลงคะแนน สั่งการโมดูลจำลองควอนตัม คำนวณค่า QBER การเพิ่มความปลอดภัยด้วยเมทริกซ์โทพลิตซ์ และถอดรหัสบัตรลงคะแนน
    \item \textbf{Data Tier (Database Storage):} ใช้ฐานข้อมูล MongoDB ร่วมกับไลบรารี Beanie ODM สำหรับจัดเก็บข้อมูลแบบ Document-oriented เพื่อบันทึกข้อมูลการเลือกตั้ง รายชื่อผู้สมัคร และประวัติเซสชัน
\end{enumerate}

\subsection{Quantum Channel Simulation}
เนื่องจากฮาร์ดแวร์คอมพิวเตอร์ควอนตัมจริงในปัจจุบันยังมีข้อจำกัดทางกายภาพค่อนข้างมาก โดยสถานะควอนตัมยังไวต่อสัญญาณรบกวนจากสภาพแวดล้อมภายนอกได้ง่าย (Physical Noise) อีกทั้งการส่งคำสั่งไปประมวลผลบนเครื่องควอนตัมจริงผ่านระบบคลาวด์ยังมีค่าใช้จ่ายและต้องรอคิวประมวลผลเป็นเวลานาน จึงได้ออกแบบระบบให้เลือกได้ระบบเพื่อให้สามารถเลือกใช้งานได้การจำลองได้สองรูปแบบ:
\begin{enumerate}
    \item \textbf{Local Quantum Simulator (Qiskit Aer):} ระบบที่จำลองการทำงานของคอมพิวเตอร์ความตัมบนเครื่องเซิร์ฟเวอร์โดยตรงผ่านไลบรารี \texttt{qiskit-aer} ให้การประมวลผลที่รวดเร็วในระดับมิลลิวินาที ไม่ต้องรอคิว และไม่มีค่าใช้จ่าย สามารถเลือกจำลองได้ทั้งช่องสัญญาณในอุดมคติที่ไร้สัญญาณรบกวน หรือสภาวะที่มีสัญญาณรบกวนแบบจำลองได้ เหมาะสำหรับการทดสอบระบบและการใช้งานทั่วไป
    \item \textbf{Cloud Quantum Hardware (IBM Quantum Runtime):} รองรับการเชื่อมต่อกับเครื่องคอมพิวเตอร์ควอนตัมจริงของ IBM ผ่าน \texttt{qiskit-ibm-runtime} โดยแปลงวงจรควอนตัมให้อยู่ในรูปของชุดคำสั่งที่ตรงกับของชิปควอนตัมจริง ทำให้สามารถศึกษาและเปรียบเทียบพฤติกรรมของโพรโทคอลบนฮาร์ดแวร์จริงได้
\end{enumerate}

% ============================================================================
% 3.2 BB84 QKD Engine and Key Distribution Workflow
% ============================================================================
\section{Quantum Key Distribution BB84 Protocol}
หัวใจสำคัญของระบบคือโมดูลการสร้างและกระจายกุญแจลับเชิงควอนตัม ซึ่งนำโพรโทคอล BB84 มาแปลงเป็นขั้นตอนวิธีเชิงวิศวกรรมคอมพิวเตอร์ โดยมีลำดับขั้นตอนการทำงาน ดังนี้:

\subsection{State Preparation by Election Server}
เมื่อผู้ลงคะแนนเสียงเริ่มต้นเซสชันใหม่ เซิร์ฟเวอร์การเลือกตั้งจะดำเนินการคำนวณจำนวนควอนตัมบิตทั้งหมด ($N$) ที่ต้องสร้างขึ้นตามขนาดกุญแจเป้าหมายและค่าเผื่อความปลอดภัย จากนั้นดำเนินการ:
\begin{enumerate}
    \item สุ่มลำดับบิตข้อมูล ความยาว $N$ บิต โดยแต่ละบิตมีค่าเป็น $0$ หรือ $1$
    \item สุ่มเลือกลำดับของแกนที่ใช้ในการส่ง ความยาว $N$ บิต โดยใช้รหัส $0$ แทนฐานแนวตรง ($+$ หรือ Z-basis) และรหัส $1$ แทนฐานแนวทแยง ($\times$ หรือ X-basis)
    \item สร้างวงจรควอนตัมขนาด $N$ คิวบิตด้วย Qiskit โดยกำหนดเกตทางควอนตัมตามตรรกะดังตารางที่~\ref{tab:qiskit_circuit_encoding}
\end{enumerate}

\begin{table}[htbp]
\centering
\renewcommand{\arraystretch}{1.35}
\setlength{\extrarowheight}{2.5pt}
\begin{tabular}{cccc}
\hline
\textbf{บิต} & \textbf{แกนการส่ง} & \textbf{เกตที่ใช้เตรียมสถานะ} & \textbf{สถานะควอนตัมผลลัพธ์} \\ \hline
\rowcolor{white}
0 & 0 ($+$: Z-basis) & ไม่มี (สถานะเริ่มต้น) & $|0\rangle$\\
\rowcolor[HTML]{E4E7EB}
1 & 0 ($+$: Z-basis) & เกต Pauli-X & $|1\rangle$\\
\rowcolor{white}
0 & 1 ($\times$: X-basis) & เกต Hadamard & $|+\rangle = \frac{1}{\sqrt{2}}(|0\rangle + |1\rangle)$\\
\rowcolor[HTML]{E4E7EB}
1 & 1 ($\times$: X-basis) & เกต Pauli-X ตามด้วย Hadamard & $|-\rangle = \frac{1}{\sqrt{2}}(|0\rangle - |1\rangle)$\\ \hline
\end{tabular}
\caption{การแปลงบิตและแกนส่งสู่เกตทางควอนตัมในวงจร Qiskit}
\label{tab:qiskit_circuit_encoding}
\end{table}

\subsection{Eavesdropping Simulation in Quantum Channel (Eve)}
เพื่อศึกษาพฤติกรรมความปลอดภัยและทดสอบกลไกการตรวจจับการโจมตี ระบบได้พัฒนาโมดูลจำลองการดักฟังในช่องสัญญาณควอนตัมแบบดักจับแล้วส่งต่อ (Intercept-and-Resend) ของผู้ดักฟัง (Eve) ดังนี้:
\begin{enumerate}
    \item \textbf{การดักสกัดโฟตอนระหว่างทาง:} หากเปิดใช้งานตัวเลือกจำลองการดักฟัง ระบบจะแทรกขั้นตอนการดักวัดของ Eve เข้าไปในวงจรควอนตัมก่อนที่โฟตอนจะเดินทางไปถึงเครื่องของผู้ลงคะแนนเสียง (Bob)
    \item \textbf{การสุ่มเลือกแกนวัดของ Eve:} เนื่องจากข้อมูลควอนตัมไม่สามารถแอบคัดลอกเก็บไว้ได้ (No-Cloning Theorem) ทำให้ Eve ไม่สามารถคัดลอกโฟตอนเก็บไว้รอฟังการประกาศแกนวัดเฉลยในภายหลังได้ Eve จึงจำเป็นต้องเดาสุ่มแกนวัด ($+$ หรือ $\times$) เพื่อวัดค่าโฟตอนแต่ละตัวในทันที:
    \begin{itemize}
        \item หาก Eve เลือกแกนแนวตรง ($+$: Z-basis): ทำการวัดค่าลงบนบิตข้อมูลของตนเองโดยตรง
        \item หาก Eve เลือกแกนแนวทแยง ($\times$: X-basis): ปรับแกนวัดด้วยเกต Hadamard ก่อนทำการวัดค่า
    \end{itemize}
    \item \textbf{การเปลี่ยนแปลงสถานะของโฟตอนและการส่งต่อ:} การที่เกิดการแอบวัดค่าจะเป็นการรบกวนและเปลี่ยนสถานะดั้งเดิมของโฟตอนให้กลายเป็นค่าที่อ่านได้ทันที ผู้ดักฟังจึงจำเป็นต้องสร้างโฟตอนตัวใหม่ตามค่าที่ตนเองวัดได้ส่งต่อไปให้ ผู้ดักฟังโฟตอนนั้นจึงไม่ใช่สถานะเดิมที่ส่งมาอีกต่อไป และหากผู้ดักฟังเดาสุ่มแกนวัดผิด ข้อมูลที่ส่งต่อไปก็จะผิดเพี้ยนตามไปด้วย ซึ่งจะทิ้งร่องรอยเป็นข้อผิดพลาดในบิตข้อมูล (Error) เฉลี่ยอย่างน้อย $25\%$ ในตำแหน่งที่ผู้ส่งและผู้รับเลือกแกนตรงกัน ทำให้ระบบสามารถตรวจพบความผิดปกตินี้ได้
\end{enumerate}

\subsection{Quantum Measurement by Voter Terminal}
เมื่อโฟตอน (หรือโฟตอนที่ผ่านการแทรกแซงจาก Eve) เดินทางมาถึงอุปกรณ์ตรวจวัดของผู้ลงคะแนนเสียง (Bob) จะดำเนินการตรวจวัดสถานะทางฟิสิกส์ดังนี้:
\begin{enumerate}
    \item \textbf{การสุ่มเลือกลำดับฐานวัด (Measurement Bases):} Bob สุ่มเลือกลำดับฐานวัด ($+$ หรือ $\times$) ของตนเองอย่างอิสระสำหรับแต่ละคิวบิต โดยไม่ทราบฐานวัดของ Alice และไม่ทราบการมีอยู่ของ Eve
    \item \textbf{การกำหนดเกตวัดค่าในวงจรควอนตัม:} นำฐานวัดที่สุ่มได้มากำหนดการวัดในวงจรควอนตัม:
    \begin{itemize}
        \item หาก Bob เลือกฐานแนวตรง ($+$: 0): ทำการวัดค่าลงบนบิตคลาสสิกโดยตรง
        \item หาก Bob เลือกฐานแนวทแยง ($\times$: 1): ใส่เกต Hadamard (\texttt{qc.h(i)}) เพื่อหมุนแกนอ้างอิงกลับสู่แกนตั้งตรงก่อนทำการวัด
    \end{itemize}
    \item \textbf{การประมวลผลและรับผลลัพธ์ดิบ:} นำวงจรควอนตัมที่สมบูรณ์ไปประมวลผลผ่าน \texttt{Sampler} ของ Qiskit Aer หรือส่งไปยังเครื่องคอมพิวเตอร์ควอนตัมจริงของ IBM Quantum เพื่อรับผลลัพธ์การวัดดิบของ Bob (Bob's Raw Read)
\end{enumerate}

\subsection{Basis Sifting via Classical Channel}
เมื่อกระบวนการวัดทางควอนตัมเสร็จสิ้น Alice (Server) และ Bob (Voter) จะทำการเปรียบเทียบฐานการวัดผ่านช่องทางสาธารณะคลาสสิก โดยฟังก์ชัน \texttt{sift\_bases(alice\_basis, bob\_basis)}:
\begin{equation}
    \text{Keep}[i] = \begin{cases}
        \text{True} & \text{ถ้า } \text{Alice\_Basis}[i] == \text{Bob\_Basis}[i] \\
        \text{False} & \text{ถ้า } \text{Alice\_Basis}[i] \neq \text{Bob\_Basis}[i]
    \end{cases}
\end{equation}
ทั้งสองฝ่ายจะเก็บเฉพาะบิตในตำแหน่งที่ $\text{Keep}[i] == \text{True}$ ไว้ เกิดเป็น \textbf{กุญแจคัดกรองของเซิร์ฟเวอร์ ($K_{\text{sift}}^{\text{Alice}}$)} และ \textbf{กุญแจคัดกรองของผู้โหวต ($K_{\text{sift}}^{\text{Bob}}$)} ซึ่งบิตในตำแหน่งที่ฐานไม่ตรงกันจะถูกทิ้งไปอย่างถาวร

\subsection{Parameter Estimation and Strict Abort Policy}
เพื่อตรวจสอบว่าช่องสัญญาณมีความปลอดภัยและไม่มีบุคคลที่สามแอบดักฟัง ระบบจะเข้าสู่ขั้นตอนการประเมินค่าพารามิเตอร์:
\begin{enumerate}
    \item ทั้งสองฝ่ายสุ่มเลือกดัชนีบิตตัวอย่างจำนวน $k$ บิตจาก Sifted Key (เรียกว่า Sample Bits) เพื่อนำมาเปิดเผยเปรียบเทียบกันบนช่องทางสาธารณะ
    \item คำนวณอัตราความผิดพลาดของบิตควอนตัมจากกลุ่มตัวอย่าง (Practical Sample QBER: $e_{\text{sample}}$):
    \begin{equation}
        e_{\text{sample}} = \frac{\sum_{j=1}^{k} (K_{\text{sift}}^{\text{Alice}}[j] \oplus K_{\text{sift}}^{\text{Bob}}[j])}{k} \times 100\%
    \end{equation}
    \item \textbf{การบังคับใช้นโยบาย Strict Abort Policy (ไม่มี Error Correction):} ระบบจะนำค่า $e_{\text{sample}}$ ไปเปรียบเทียบกับเพดานความปลอดภัย $11.0\%$~\cite{shor_preskill_2000}:
    \begin{itemize}
        \item หาก $e_{\text{sample}} > 11.0\%$ หรือในโหมดการจำลองที่ไร้สัญญาณรบกวนตรวจพบบิตผิดพลาดแม้แต่บิตเดียว ($e > 0\%$): เซิร์ฟเวอร์จะเปลี่ยนสถานะเซสชันเป็น \texttt{ABORTED} สั่งยกเลิกการลงคะแนนทันที และทำลายกุญแจทั้งหมดในหน่วยความจำทิ้ง โดย\textbf{ไม่มีการใช้ขั้นตอนวิธีแก้ไขข้อผิดพลาด (Key Reconciliation)} เพื่อป้องกันการรั่วไหลของข้อมูลผ่าน Parity Bits
        \item หาก $e_{\text{sample}} \le 11.0\%$: กุญแจได้รับการรับรองความปลอดภัย บิตตัวอย่างที่ใช้ตรวจสอบจะถูกตัดทิ้ง และบิตที่เหลือรอดจะถูกส่งต่อไปยังขั้นตอนการขยายความลับ
    \end{itemize}
\end{enumerate}

\subsection{Privacy Amplification via Toeplitz Hashing}
เพื่อกำจัดข้อมูลที่ผู้ดักฟังอาจแอบล่วงรู้ไปในระหว่างที่ค่า QBER ยังไม่เกินเกณฑ์ความปลอดภัย ระบบจะนำกุญแจดิบส่วนที่เหลือ (ความยาว $n_{\text{raw}}$ บิต) มาผ่านขั้นตอนการขยายความลับ (Privacy Amplification):
\begin{enumerate}
    \item เซิร์ฟเวอร์สุ่มสตริงตัวเลขฐานสองสาธารณะเรียกว่า \textbf{Seed ($S$)} ความยาว $n_{\text{raw}} + m - 1$ บิต โดยที่ $m$ คือขนาดกุญแจเป้าหมาย (Target Key Length) และส่ง $S$ ให้แก่ผู้โหวตผ่านช่องทางคลาสสิก
    \item ทั้ง Server (Alice) และ Voter (Bob) สร้างเมทริกซ์โทพลิตซ์ $T$ ขนาด $m \times n_{\text{raw}}$ จาก Seed เดียวกัน โดยสมาชิกในแนวทแยงขนานกันมีค่าเท่ากัน:
    \begin{equation}
        T_{i, j} = S[n_{\text{raw}} - 1 + i - j]
    \end{equation}
    \item ทำการคูณเมทริกซ์ $T$ เข้ากับเวกเตอร์กุญแจดิบ $X$ ภายใต้เลขคณิตมอดุโลสอง (การดำเนินการระดับบิต AND และ XOR) ในฟังก์ชัน \texttt{bitwise\_dot\_product}:
    \begin{equation}
        K_{\text{final}}[i] = \bigoplus_{j=0}^{n_{\text{raw}}-1} (T_{i, j} \land X[j]) \quad \text{สำหรับ } i = 0, 1, \dots, m-1
    \end{equation}
\end{enumerate}
ผลลัพธ์ที่ได้คือ \textbf{กุญแจลับบริสุทธิ์สุดท้าย ($K_{\text{final}}$)} ความยาว $m$ บิต ที่สมบูรณ์ ปลอดภัย และตรงกันทุกประการระหว่าง Server และ Voter

\subsection{Finite-Key Safety Margin Calculation}
เนื่องจากจำนวนบิตที่ผ่านการคัดกรองฐานวัดมีความผันผวนตามการแจกแจงแบบทวินาม หากระบบส่งโฟตอนเพียง $4m$ ตัว จะเกิดความเสี่ยงสูงที่บิตจะไม่เพียงพอต่อการหักบิตตัวอย่างและบีบอัดกุญแจ โมดูล \texttt{service.py} จึงนำสมการค่าเผื่อความปลอดภัยมาใช้คำนวณจำนวนควอนตัมบิตที่ต้องส่งจริง:
\begin{equation}
    N = \left\lceil (4 + \delta)(m + s) \right\rceil
\end{equation}
โดยในโค้ดของระบบได้กำหนดค่าเริ่มต้นของพารามิเตอร์ความปลอดภัย $s = 16$ บิต และกำหนดค่าเผื่อความปลอดภัยที่เหมาะสม $\delta = 0.8$ ซึ่งรับประกันอัตราความสำเร็จในการสร้างกุญแจสูงกว่า $97\%$ ในทุกเซสชัน

% ============================================================================
% 3.3 One-Time Pad Ballot Casting and Tallying
% ============================================================================
\section{One-Time Pad Ballot Casting and Tallying}
เมื่อกระบวนการ QKD สร้างกุญแจลับบริสุทธิ์ $K_{\text{final}}$ สำเร็จแล้ว ระบบจะนำกุญแจนี้ไปใช้ในการปกป้องความลับของบัตรลงคะแนนเสียงด้วยขั้นตอนวิธีการเข้ารหัสที่ปลอดภัยที่สุดในโลก คือ One-Time Pad (OTP)

\subsection{Ballot Representation and Binary Encoding}
ในการเลือกตั้ง ผู้สมัครแต่ละรายจะได้รับหมายเลขประจำตัวหรือดัชนี (Candidate Index) ระบบจะแปลงตัวเลือกของผู้มีสิทธิ์ลงคะแนนให้เป็นลำดับบิตข้อความต้นฉบับ (Plaintext Ballot: $M$) ที่มีความยาวเท่ากับขนาดของกุญแจลับ $K_{\text{final}}$ พอดี ($|M| = |K_{\text{final}}| = m$ บิต) เช่น การใช้รหัสฐานสองแทนหมายเลขผู้สมัคร หรือการใช้เวกเตอร์แบบ One-Hot Encoding เพื่อให้สามารถดำเนินการระดับบิตได้อย่างสมบูรณ์

\subsection{Client-Side OTP Encryption (Bob)}
ผู้ลงคะแนนเสียง (Bob) จะทำการเข้ารหัสบัตรลงคะแนนเสียงบนเว็บเบราว์เซอร์ของตนเอง โดยไม่มีการส่งตัวเลือกที่ยังไม่ได้เข้ารหัสออกสู่เครือข่าย กระบวนการเข้ารหัสใช้ตัวดำเนินการ XOR แบบบิตต่อบิต:
\begin{equation}
    C = M \oplus K_{\text{final}} \implies C[i] = M[i] \oplus K_{\text{final}}[i] \quad \text{สำหรับทุก } i
\end{equation}
ข้อความไซเฟอร์ $C$ (Encrypted Vote) ที่ได้จะมีลักษณะเป็นสตริงสุ่มโดยสมบูรณ์ จากนั้น Voter Client จะส่ง $C$ พร้อมด้วยรหัสเซสชัน (\texttt{sessionId}) และรหัสผู้โหวต (\texttt{voterId}) ไปยังเซิร์ฟเวอร์ผ่านคำขอ \texttt{POST /vote/cast-ballot} บนช่องทางเครือข่ายสาธารณะ

\subsection{Server-Side Decryption, Secure Tallying, and Key Destruction (Alice)}
เมื่อเซิร์ฟเวอร์การเลือกตั้ง (Alice) ได้รับข้อความไซเฟอร์ $C$ จะดำเนินการตามลำดับขั้นดังนี้:
\begin{enumerate}
    \item ดึงกุญแจลับ $K_{\text{final}}$ ประจำเซสชันนั้นที่ Alice เก็บไว้ในหน่วยความจำ
    \item ทำการถอดรหัสบัตรลงคะแนนเสียงด้วยตัวดำเนินการ XOR ร่วมกับกุญแจชุดเดิม:
    \begin{equation}
        M = C \oplus K_{\text{final}}
    \end{equation}
    ตามสมบัติการหักล้างตัวเองของ XOR ทำให้ได้ข้อมูลการลงคะแนน $M$ ดั้งเดิมกลับคืนมาอย่างถูกต้อง $100\%$
    \item ปรับปรุงผลคะแนนสะสมของผู้สมัครที่ได้รับเลือกในฐานข้อมูล (\texttt{candidates.votes += 1})
    \item \textbf{การทำลายกุญแจทิ้งทันที (Strict Key Destruction):} เซิร์ฟเวอร์จะทำการลบกุญแจ $K_{\text{final}}$ ออกจากหน่วยความจำและเปลี่ยนสถานะเซสชันในฐานข้อมูลเป็น \texttt{VOTE\_CAST} โดยไม่มีการบันทึกกุญแจลับลงในดิสก์ถาวร เพื่อให้เป็นไปตามกฎเหล็กของ One-Time Pad ว่ากุญแจแต่ละชุดต้องใช้เพียงครั้งเดียวแล้วทำลายทิ้งทันที (Never Reused) ป้องกันช่องโหว่ Two-Time Pad โดยเด็ดขาด
\end{enumerate}

% ============================================================================
% 3.4 System Implementation Details
% ============================================================================
\section{System Implementation Details}
ในส่วนนี้จะกล่าวถึงรายละเอียดเชิงวิศวกรรม การเลือกใช้เทคโนโลยี และโครงสร้างสถาปัตยกรรมภายในของซอฟต์แวร์ที่พัฒนาขึ้นจริง

\subsection{Frontend Architecture and Multi-Perspective Interface}
ส่วนติดต่อผู้ใช้งาน (Frontend) พัฒนาขึ้นในรูปแบบ Single Page Application (SPA) ด้วย React เวอร์ชัน 18 ร่วมกับภาษา TypeScript และใช้ Vite เป็นเครื่องมือสำหรับการคอมไพล์และประกอบโปรแกรม (Build Tool) โครงสร้างซอฟต์แวร์ถูกจัดระเบียบตามขอบเขตความรับผิดชอบ ได้แก่ หน้าจอหลัก (\texttt{pages/}), ชิ้นส่วนคอมโพเนนต์ที่นำกลับมาใช้ซ้ำ (\texttt{components/}), บริการเชื่อมต่อเครือข่าย API (\texttt{services/}), และระบบจัดการสถานะข้อมูลระดับส่วนกลาง (\texttt{context/}) เพื่อควบคุมวงจรชีวิตของการเลือกตั้งตั้งแต่การยืนยันตัวตนจนถึงการรับใบเสร็จยืนยันการลงคะแนน

ในการจำลองระบบการสื่อสารควอนตัม จุดท้าทายสำคัญคือ \textbf{"ระดับการเข้าถึงสารสนเทศ (Information Horizon)"} ของผู้ใช้งานแต่ละกลุ่มในสถานการณ์จริงมีความแตกต่างกันอย่างมาก การเปิดเผยข้อมูลภายในช่องสัญญาณทั้งหมดให้แก่ผู้โหวตทั่วไปจะทำให้ระบบไม่สะท้อนความเป็นจริงทางกายภาพ ในทางกลับกัน หากปิดบังกระบวนการทางฟิสิกส์ทั้งหมด ก็จะไม่สามารถใช้เป็นเครื่องมือสำหรับศึกษาและตรวจสอบความโปร่งใสได้ ระบบจึงได้ออกแบบหน้า \textbf{KeyGenerationPage} ให้รองรับการแสดงผล \textbf{3 มุมมอง (Multi-Perspective Visualization)} เพื่อตอบโจทย์กลุ่มผู้ใช้งานที่หลากหลาย:

\subsubsection{1. Voter Terminal View (มุมมองผู้ลงคะแนนเสียง)}
มุมมองนี้จำลองสภาพแวดล้อมที่ผู้มีสิทธิ์เลือกตั้ง (Bob) จะได้เห็นจริงบนหน้าจอเครื่องลงคะแนนเสียง (Voting Terminal) ในคูหาเลือกตั้ง โดยยึดหลัก \textbf{การจำกัดขอบเขตข้อมูลเท่าที่จำเป็น (Principle of Least Privilege \& Realistic Emulation)}:
\begin{itemize}
    \item \textbf{ข้อมูลที่แสดงผล:} แสดงเฉพาะผลลัพธ์ที่ตรวจวัดได้จากอุปกรณ์ตรวจจับโฟตอน (Photon Detector) ภายในเครื่องของ Bob เท่านั้น ได้แก่ ลำดับฐานวัดที่ Bob สุ่มเลือก และบิตผลลัพธ์ที่อ่านได้ พร้อมทั้งแสดงผลการแจ้งเตือนสาธารณะจาก Server ว่าโฟตอนตำแหน่งใดมีฐานวัดตรงกัน (Keep) หรือไม่ตรงกัน (Discard) และแสดงดัชนีของบิตตัวอย่างที่ถูกสุ่มเลือกไปใช้ตรวจสอบความปลอดภัย
    \item \textbf{การบ่งชี้ความปลอดภัย:} นำเสนอผ่านตัวชี้วัดสถานะที่เข้าใจง่าย ได้แก่ แถบความคืบหน้าของการรับส่งโฟตอน และเกจวัดความปลอดภัย หากค่าความผิดพลาดของบิตตัวอย่างไม่เกินเกณฑ์ความปลอดภัย $11.0\%$ ระบบจะขึ้นสถานะสีเขียวรับรองความปลอดภัยและปลดล็อกปุ่มดำเนินการลงคะแนนเสียง แต่หากตรวจพบความผิดปกติ ระบบจะแสดงข้อความเตือนและสั่งยกเลิกเซสชันทันที
    \item \textbf{ข้อมูลที่ถูกซ่อนไว้:} บิตดั้งเดิมและฐานวัดของ Alice จะไม่ถูกเปิดเผยให้เห็น ตลอดจนไม่มีการแสดงตัวตนของผู้ดักฟัง (Eve) เพื่อสะท้อนความเป็นจริงว่าในระบบกายภาพ ผู้ใช้งานปลายทางจะไม่สามารถล่วงรู้กระบวนการภายในช่องสัญญาณภายนอกได้เลย
\end{itemize}

\subsubsection{2. Practical Channel View (มุมมองช่องสัญญาณเชิงปฏิบัติการ)}
มุมมองนี้ออกแบบมาสำหรับผู้ดูแลระบบ ผู้สังเกตการณ์การเลือกตั้ง และวิศวกรความปลอดภัยเครือข่าย (Security Auditor) เพื่อตรวจสอบความสมบูรณ์และเสถียรภาพของการเชื่อมต่อระหว่างสองโหนดสื่อสาร:
\begin{itemize}
    \item \textbf{ข้อมูลที่แสดงผล:} แสดงภาพจำลองช่องสัญญาณควอนตัมที่เชื่อมต่อระหว่าง Election Server (Sender) และ Voter Terminal (Receiver) โดยจับคู่เปรียบเทียบฐานวัดของทั้งสองฝ่ายแบบคู่ขนาน ทำให้มองเห็นกระบวนการคัดกรองฐานวัด (Basis Sifting) ได้อย่างชัดเจน ตำแหน่งที่ฐานตรงกันจะถูกไฮไลต์ด้วยสีเขียว ส่วนตำแหน่งที่ฐานต่างกันจะถูกหรี่แสงลง
    \item \textbf{การตรวจจับการดักฟังแบบซ่อนเร้น (Hidden Interceptor):} ในมุมมองนี้ ผู้ดักฟัง (Eve) จะถูกจำลองในฐานะบุคคลที่สามที่แอบแทรกแซงทางกายภาพอยู่ในสายส่ง ซึ่งระบบจะไม่สามารถมองเห็นตัวตน หน้าจอ หรือฐานการวัดของ Eve ได้โดยตรง แต่ระบบจะ \textbf{"อนุมานการมีอยู่ของ Eve ผ่านผลกระทบทางสถิติ"} โดยมีมิเตอร์วัดค่า Practical Sample QBER เปรียบเทียบกับเส้นประสีแดงของเพดานความปลอดภัย $11.0\%$ หากมีผู้ดักฟัง ค่า QBER จะกระโดดสูงขึ้นสู่ระดับ $25.0\%$ อย่างเด่นชัด ทำให้ผู้ดูแลระบบตรวจพบการบุกรุกได้อย่างแม่นยำ
    \item \textbf{ตัวชี้วัดเชิงปฏิบัติการ:} แสดงอัตราส่วนการบีบอัดกุญแจ (Compression Ratio) ของกระบวนการ Privacy Amplification และปริมาณข้อมูลสถิติของช่องสัญญาณ
\end{itemize}

\subsubsection{3. System View (มุมมองการศึกษาเชิงลึกแบบ Omniscient)}
มุมมองนี้ถูกสร้างขึ้นเพื่อวัตถุประสงค์ทางการศึกษาและการวิจัยทางวิชาการ (Educational \& Scientific Ground Truth) โดยทำหน้าที่เป็น \textbf{"มุมมองของผู้รู้แจ้งทุกสิ่ง (Omniscient Observer)"} ที่เปิดเผยกระบวนการทางควอนตัมทั้งหมดที่เกิดขึ้นในสายส่งแบบโปร่งใส $100\%$:
\begin{itemize}
    \item \textbf{บันทึกการส่งผ่านโฟตอนแบบครบวงจร (Full Physical Trace):} แสดงตารางสถานะของทุกโฟตอนตั้งแต่ต้นทางจนถึงปลายทาง ประกอบด้วย:
    \begin{enumerate}
        \item บิตข้อมูลดิบ ฐานวัด และมุมโพลาไรเซชันของโฟตอนที่ Alice ส่งออก
        \item การแทรกแซงของ Eve: ฐานที่ Eve สุ่มวัด ค่าที่ Eve อ่านได้ และสถานะโฟตอนตัวใหม่ที่ Eve สร้างขึ้นเพื่อส่งต่อ
        \item ฐานวัดและค่าที่ Bob อ่านได้จากอุปกรณ์ตรวจจับ
    \end{enumerate}
    \item \textbf{การชี้เป้าความผิดพลาดทางฟิสิกส์ (Collapse Error Highlighting):} ในตำแหน่งที่ Alice และ Bob ใช้ฐานตรงกัน แต่ Eve เข้ามาดักวัดด้วยฐานที่ผิด ระบบจะทำการไฮไลต์แถวข้อมูลด้วยสีแดง เพื่อชี้ให้เห็นอย่างเป็นรูปธรรมว่า การยุบตัวของฟังก์ชันคลื่นตามกฎของบอร์นทำให้บิตของ Bob ผิดพลาดไปจาก Alice อย่างไร ซึ่งเป็นหลักฐานเชิงประจักษ์ของรอยนิ้วมือ $25\%$ QBER
    \item \textbf{การเปรียบเทียบค่าความผิดพลาดจริงกับกลุ่มตัวอย่าง:} แสดงการเปรียบเทียบระหว่างค่าความผิดพลาดจริงของทั้งระบบ (Ground Truth QBER: $e_{\text{system}}$) กับค่าที่ประมาณได้จากกลุ่มตัวอย่างสาธารณะ (Sample QBER: $e_{\text{sample}}$) เพื่อให้นักศึกษาสามารถตรวจสอบความถูกต้องของการอนุมานทางสถิติได้
    \item \textbf{การตรวจสอบกระบวนการขยายความลับ:} แสดงเมล็ดสุ่มสาธารณะ (PA Seed), โครงสร้างของเมทริกซ์โทพลิตซ์ $T$, และเวกเตอร์กุญแจดิบที่ถูกนำมาบีบอัด เพื่อสาธิตให้เห็นว่าข้อมูลส่วนที่อาจรั่วไหลไปยัง Eve ถูกตัดทิ้งไปจนเหลือศูนย์อย่างสมบูรณ์
\end{itemize}

ตารางที่~\ref{tab:view_modes_comparison} สรุปการเปรียบเทียบขอบเขตและคุณลักษณะของมุมมองการแสดงผลทั้ง 3 รูปแบบในระบบ

\begin{table}[htbp]
\centering
\renewcommand{\arraystretch}{1.35}
\setlength{\extrarowheight}{2pt}
\caption{การเปรียบเทียบขอบเขตการเข้าถึงข้อมูลและคุณลักษณะของมุมมองทั้ง 3 รูปแบบ}
\label{tab:view_modes_comparison}
\begin{tabularx}{\textwidth}{>{\raggedright\arraybackslash}p{2.6cm} >{\raggedright\arraybackslash}X >{\raggedright\arraybackslash}X >{\raggedright\arraybackslash}X}
\hline
\textbf{หัวข้อ} & \textbf{Voter View} & \textbf{Practical View} & \textbf{System View} \\ \hline
\rowcolor{white}
\textbf{กลุ่มเป้าหมาย} & ประชาชนผู้มาใช้สิทธิ์โหวต & ผู้ดูแลระบบและผู้ตรวจสอบ & นักวิจัยและนักศึกษา \\
\rowcolor[HTML]{E4E7EB}
\textbf{ขอบเขตข้อมูล} & เฉพาะที่เครื่องตรวจวัดของตน & ช่องสัญญาณ Alice $\to$ Bob & ครบวงจรทุกขั้นตอน (Omniscient) \\
\rowcolor{white}
\textbf{การมองเห็น Eve} & มองไม่เห็นโดยเด่นชัด & เห็นเป็นรอยนิ้วมือสถิติ QBER & เห็นฐาน บิต และวงจรของ Eve ทั้งหมด \\
\rowcolor[HTML]{E4E7EB}
\textbf{การแสดงค่า QBER} & สถานะความปลอดภัย (Safe/Abort) & Practical Sample QBER & ทั้ง Sample QBER และ Ground Truth \\
\rowcolor{white}
\textbf{จุดประสงค์หลัก} & ใช้งานง่าย ลดความสับสนในการโหวต & เฝ้าระวังและตรวจจับการโจมตี & ศึกษาฟิสิกส์ควอนตัมเชิงลึก \\ \hline
\end{tabularx}
\end{table}

\subsection{Backend Architecture and REST API Specification}
ส่วนเซิร์ฟเวอร์ประมวลผล (Backend) พัฒนาด้วยภาษา Python บนเฟรมเวิร์ค FastAPI มีโครงสร้างการทำงานแบบ Asynchronous ที่มีประสิทธิภาพสูง แบ่งโมดูลการทำงานภายในออกเป็น:
\begin{itemize}
    \item \texttt{modules/qkd.py}: รับผิดชอบตรรกะการจำลองวงจรควอนตัมของโพรโทคอล BB84 ผ่าน Qiskit การสุ่มบิต การวัด และการดักฟัง
    \item \texttt{modules/privacy\_amplification.py}: รับผิดชอบการสร้างเมทริกซ์โทพลิตซ์จาก Seed และการคำนวณการคูณเมทริกซ์ระดับบิต
    \item \texttt{modules/service.py}: ควบคุมตรรกะทางธุรกิจ การตรวจสอบเงื่อนไขความปลอดภัย การจัดการเซสชัน และการเข้ารหัส/ถอดรหัส OTP
    \item \texttt{modules/router.py}: ให้บริการ RESTful API สำหรับการเชื่อมต่อระหว่าง Frontend และ Backend
\end{itemize}

ตารางที่~\ref{tab:api_endpoints} แสดงรายละเอียดข้อกำหนดของ REST API Endpoints หลักในระบบ

\begin{table}[htbp]
\centering
\renewcommand{\arraystretch}{1.3}
\setlength{\extrarowheight}{2pt}
\caption{ข้อกำหนดรายละเอียดของ REST API Endpoints ในระบบจำลองการลงคะแนนเสียง}
\label{tab:api_endpoints}
\begin{tabularx}{\textwidth}{p{1.4cm}p{3.8cm}p{3.2cm}X}
\hline
\textbf{Method} & \textbf{Endpoint Path} & \textbf{Request Payload} & \textbf{คำอธิบายหน้าที่การทำงาน} \\ \hline
\rowcolor{white}
\texttt{POST} & \texttt{/vote/generate-key} & \texttt{GenerateKeyInput} & เริ่มต้นเซสชัน QKD สร้างวงจร BB84 คัดกรองฐาน ตรวจสอบ QBER และคืนค่ากุญแจให้แก่ Bob \\
\rowcolor[HTML]{E4E7EB}
\texttt{POST} & \texttt{/vote/cast-ballot} & \texttt{CastBallotInput} & รับบัตรลงคะแนนที่เข้ารหัสด้วย OTP ($C$) มาถอดรหัส นับคะแนน และทำลายกุญแจ \\
\rowcolor{white}
\texttt{POST} & \texttt{/vote/abort-session} & \texttt{AbortSessionInput} & สั่งยกเลิกเซสชันและทำลายกุญแจเมื่อผู้ใช้งานหรือระบบตรวจพบความผิดปกติ \\
\rowcolor[HTML]{E4E7EB}
\texttt{GET} & \texttt{/vote/candidates/\{code\}} & None & ดึงรายชื่อผู้สมัครและพรรคการเมืองในการเลือกตั้งที่ระบุ \\
\rowcolor{white}
\texttt{GET} & \texttt{/vote/election-results/\{code\}} & None & ดึงรายงานผลคะแนนเสียงสะสมของการเลือกตั้ง \\
\rowcolor[HTML]{E4E7EB}
\texttt{GET} & \texttt{/vote/logs/\{code\}} & None & ดึงประวัติ Quantum Audit Log เพื่อการตรวจสอบย้อนหลัง \\ \hline
\end{tabularx}
\end{table}

\subsection{Database Schema and Quantum Audit Log}
การจัดเก็บข้อมูลใช้ฐานข้อมูล MongoDB โดยมีโครงสร้างเอกสาร (Document Schemas) ที่กำหนดผ่าน Beanie ODM 2 คอลเลกชันหลัก:
\begin{enumerate}
    \item \textbf{ElectionSchema:} จัดเก็บข้อมูลหลักของการเลือกตั้ง ประกอบด้วย รหัสการเลือกตั้ง (\texttt{election\_code}), รายชื่อผู้สมัครและคะแนนเสียงสะสม (\texttt{candidates}: รายชื่อวัตถุ \texttt{Candidate} ที่มีฟิลด์ \texttt{candidate\_name}, \texttt{candidate\_party}, \texttt{votes}), และประวัติเซสชันทั้งหมด (\texttt{sessions})
    \item \textbf{VotingSession:} บันทึกประวัติและตัวชี้วัดทางควอนตัมของแต่ละเซสชันการลงคะแนน (Quantum Audit Log) เพื่อให้ผู้มีสิทธิ์โหวตและผู้สังเกตการณ์สามารถตรวจสอบความโปร่งใสของกระบวนการทางฟิสิกส์ได้ ข้อมูลที่บันทึกประกอบด้วย:
    \begin{itemize}
        \item รหัสเซสชัน (\texttt{session\_id}) และรหัสผู้โหวต (\texttt{voter\_id})
        \item สถานะของเซสชัน (\texttt{status}): เช่น \texttt{KEY\_GENERATED}, \texttt{VOTE\_CAST}, หรือ \texttt{ABORTED}
        \item ค่าทางสถิติของควอนตัม: อัตราความผิดพลาดกลุ่มตัวอย่าง (\texttt{qber\_practical}), ค่าความผิดพลาดจริงของระบบ (\texttt{qber\_system}), และเกณฑ์ความปลอดภัย (\texttt{threshold\_percent} = $11.0\%$)
        \item ข้อมูลการขยายความลับ: เมล็ดสุ่ม (\texttt{pa\_seed}), อัตราการบีบอัด (\texttt{pa\_compression\_ratio}), และขอบเขตข้อมูลที่อาจรั่วไหล (\texttt{eve\_info\_bound})
        \item บัตรลงคะแนนที่ถูกเข้ารหัส (\texttt{encrypted\_vote})
    \end{itemize}
\end{enumerate}

\noindent \textbf{การรักษาความเป็นส่วนตัวของผู้ลงคะแนนเสียง (Voter Secrecy):}\\
ในระบบบันทึก Audit Log ข้อมูลบัตรลงคะแนนจะถูกจัดเก็บในรูปของข้อความไซเฟอร์ (\texttt{encrypted\_vote}) เท่านั้น และกุญแจลับ $K_{\text{final}}$ จะไม่ถูกบันทึกเก็บไว้ ทำให้แม้บุคคลภายนอกหรือผู้ดูแลฐานข้อมูลจะสามารถเข้าถึงตาราง Audit Log เพื่อตรวจสอบค่าทางสถิติ QBER ได้ แต่จะไม่สามารถถอดรหัสเพื่อล่วงรู้ได้ว่าผู้มีสิทธิ์เลือกตั้งแต่ละคนลงคะแนนให้แก่ผู้สมัครรายใด เป็นการรักษาความลับของผู้ลงคะแนนเสียงตามหลักการประชาธิปไตยอย่างสมบูรณ์

% ============================================================================
% 3.5 End-to-End System Execution Sequence
% ============================================================================
\section{End-to-End System Execution Sequence}
เพื่อให้เห็นภาพรวมการทำงานประสานกันระหว่างทุกส่วนประกอบ ตารางที่~\ref{tab:end_to_end_flow} สรุปขั้นตอนการทำงานตั้งแต่เริ่มต้นจนเสร็จสิ้นกระบวนการลงคะแนนเสียง 1 เซสชัน

\begin{table}[htbp]
\centering
\renewcommand{\arraystretch}{1.3}
\setlength{\extrarowheight}{2pt}
\caption{ลำดับขั้นตอนการทำงานรวมตั้งแต่เริ่มต้นสร้างกุญแจจนถึงการบันทึกคะแนนเสียง (End-to-End Workflow)}
\label{tab:end_to_end_flow}
\begin{tabularx}{\textwidth}{c p{2.2cm} p{2.2cm} X}
\hline
\textbf{ขั้นตอน} & \textbf{ผู้ดำเนินการ} & \textbf{ช่องทางการสื่อสาร} & \textbf{รายละเอียดกิจกรรมและผลลัพธ์} \\ \hline
\rowcolor{white}
1 & Voter (Bob) & Classical (HTTP) & เข้าสู่ระบบ ระบุ \texttt{voterId} และส่งคำขอสร้างกุญแจ (\texttt{POST /vote/generate-key}) \\
\rowcolor[HTML]{E4E7EB}
2 & Server (Alice) & Quantum Channel & สุ่มบิตและฐานวัด สร้างวงจร Qiskit และส่งคิวบิตจำลอง $N$ โฟตอนให้ Bob \\
\rowcolor{white}
3 & Voter (Bob) & Quantum Simulator & สุ่มฐานวัดของตนเอง ทำการวัดค่าคิวบิตที่ได้รับ \\
\rowcolor[HTML]{E4E7EB}
4 & Server \& Voter & Classical Channel & แลกเปลี่ยนฐานวัดเพื่อทำ Basis Sifting และสละบิตตัวอย่างตรวจสอบค่า QBER \\
\rowcolor{white}
5 & Server (Alice) & Internal Logic & ตรวจสอบเงื่อนไขความปลอดภัย:
\begin{itemize}
    \item หาก $\text{QBER} > 11.0\%$: สั่ง \texttt{ABORT} ทำลายกุญแจทันที
    \item หาก $\text{QBER} \le 11.0\%$: สุ่ม Seed และทำ Privacy Amplification
\end{itemize} \\
\rowcolor[HTML]{E4E7EB}
6 & Server $\to$ Voter & Classical (HTTP) & ส่งมอบผลลัพธ์กุญแจลับบริสุทธิ์ $K_{\text{final}}$ และ Seed ให้แก่ Voter Client \\
\rowcolor{white}
7 & Voter (Bob) & Local Browser & นำบัตรลงคะแนน $M$ มาทำ XOR กับ $K_{\text{final}}$ ได้เป็น $C = M \oplus K_{\text{final}}$ \\
\rowcolor[HTML]{E4E7EB}
8 & Voter $\to$ Server & Classical (HTTP) & ส่งข้อความไซเฟอร์ $C$ ไปยังเซิร์ฟเวอร์ (\texttt{POST /vote/cast-ballot}) \\
\rowcolor{white}
9 & Server (Alice) & Classical Server & นำ $C$ มา XOR กับ $K_{\text{final}}$ เพื่อถอดรหัสเป็น $M$, ปรับปรุงผลคะแนนใน MongoDB, และทำลายกุญแจทิ้งทันที \\
\rowcolor[HTML]{E4E7EB}
10 & Voter \& Public & Classical (HTTP) & ตรวจสอบใบเสร็จการโหวตและประวัติเซสชันในหน้า Quantum Audit Log (\texttt{/logs}) \\ \hline
\end{tabularx}
\end{table}

สถาปัตยกรรมและขั้นตอนการทำงานที่นำเสนอในบทนี้ แสดงให้เห็นถึงการบูรณาการหลักการฟิสิกส์ควอนตัม วิทยาการรหัสลับ และวิศวกรรมซอฟต์แวร์สมัยใหม่เข้าด้วยกันอย่างเป็นรูปธรรม ซึ่งสร้างรากฐานที่มั่นคงสำหรับการประเมินผลและการทดลองในบทถัดไป
